\par\medskip\noindent\textbf{Cómo leer la salida.}
{\small\setlength{\tabcolsep}{4pt}\renewcommand{\arraystretch}{1.15}
\begin{xltabular}{\linewidth}{@{}>{\raggedright\arraybackslash}p{0.24\linewidth}>{\raggedleft\arraybackslash}p{0.15\linewidth}>{\raggedright\arraybackslash}X@{}}
\toprule Indicador & Valor & Interpretación \\ \midrule \endfirsthead
\toprule Indicador & Valor & Interpretación \\ \midrule \endhead
\bottomrule \endlastfoot
\texttt{T}, media mensual & 819; 0,0070 & Rendimiento medio mensual del factor momentum (UMD): 8,37\,\% anual. \\
Desv. est. anual & 13,49\,\% & Volatilidad anual del factor. \\
$t$ clásico / Newey-West & 5,13 / 5,06 & La media es muy significativa; corregir por autocorrelación casi no cambia el resultado. \\
Razón de Sharpe anual & 0,620 & Rendimiento por unidad de riesgo: más alto que el del mercado en el mismo período. \\
Asimetría & $-1{,}378$ & Fuertemente negativa: momentum gana poco a poco y pierde mucho de golpe. \\
Peor mes & 2009-04 ($-34{,}58$\,\%) & El desplome de momentum en el rebote de los mercados tras la crisis financiera. \\
Caída máxima acumulada & $-57{,}6$\,\% (fondo 2009-09) & Pérdida desde el máximo previo hasta el mínimo: un inversionista perdió más de la mitad. \\
1949--1999 & 10,64\,\%; $t=7{,}01$ & En la primera mitad la prima es grande y muy significativa. \\
2000--2017 & 1,67\,\%; $t=0{,}36$ & Desde 2000 la prima casi desaparece y no es significativa. \\
\end{xltabular}}
